% GENERATED FILE. Edit canonical lessons, then run npm run generate:books.
% canonical-source-hash: fnv1a64:ca2b0d32d5b1501a
% canonical-lessons: ZH-C322-dangnian, ZH-C322-zaonian, ZH-C322-wannian, ZH-C322-yisheng, ZH-C322-huainian

\chapter{Back Then}
\label{ch:zh-back-then}
\hlchaptermodality{322}

\begin{chapteropening}
\textbf{By the end of this chapter:} \emph{I can talk about back then, someone's early and later years, a whole life and missing the past.}

Complete the last lesson of chapter 322: I can talk about back then, someone's early and later years, a whole life and missing the past.
\end{chapteropening}

\section[\texorpdfstring{\zh{当年} (dāngnián)}{当年 (dāngnián)}]{\zh{当年} (dāngnián) — back then; in those years}

\label{lesson:ZH-C322-dangnian}



\begin{quote}
Before the new one: say the Mandarin for age-old, then the Mandarin for a script for a language.
\end{quote}

\subsection*{You'll want to know: \zh{当年}}
\textbf{\zh{当年}} — \emph{dāngnián} — \textquotedblleft{}back then; in those years\textquotedblright{}.

\begin{quote}
\textbf{\zh{当年他常常给我写信}} \emph{dāngnián tā chángcháng gěi wǒ xiě xìn} — Back then he often sent me letters.
\end{quote}

\begin{cousinweb}
\textbf{\zh{当年}} (\emph{dāngnián}) is built from \zh{当} 'to face, that very' and \zh{年} 'a year', both inherited from Middle Chinese, \emph{tang} and \emph{nen} (Guangyun, in Baxter's transcription), which glosses the first 'to face' and 'to match' and says that when the grain ripens it is called a year.
\end{cousinweb}

\begin{grammarlens}[title={What you've built}]
One more word for telling what happened.
\end{grammarlens}

\subsection*{Your turn}
\begin{itemize}
\raggedright
  \item \emph{Say it:} \emph{dāngnián}
  \item \emph{Say it:} \emph{dāngnián}, once more
  \item \emph{Say it:} \emph{dāngnián tā chángcháng gěi wǒ xiě xìn} — Back then he often sent me letters.
  \item \emph{Recall:} say the Mandarin for age-old, then the Mandarin for a script for a language, then say \emph{dāngnián} again
\end{itemize}

\begin{tcolorbox}[breakable,colback=teal!4,colframe=teal!35!black,title={Before you move on}]
What does \zh{当年} mean? (\textquotedblleft{}Back then; in those years\textquotedblright{}.) Say it once more.
\end{tcolorbox}
\section[\texorpdfstring{\zh{早年} (zǎonián)}{早年 (zǎonián)}]{\zh{早年} (zǎonián) — one's early years; long ago}

\label{lesson:ZH-C322-zaonian}



\begin{quote}
Before the new one: say the Mandarin for a script for a language, then the Mandarin for in those years long ago.
\end{quote}

\subsection*{You'll want to know: \zh{早年}}
\textbf{\zh{早年}} — \emph{zǎonián} — \textquotedblleft{}one's early years; long ago\textquotedblright{}.

\begin{quote}
\textbf{\zh{他早年在外国工作过}} \emph{tā zǎonián zài wàiguó gōngzuò guo} — In his early years he worked abroad.
\end{quote}

\begin{cousinweb}
\textbf{\zh{早年}} (\emph{zǎonián}) is built from \zh{早} 'early' and \zh{年} 'a year', both inherited from Middle Chinese, \emph{tsawX} and \emph{nen} (Guangyun, in Baxter's transcription), which glosses the first 'morning'.
\end{cousinweb}

\begin{grammarlens}[title={What you've built}]
One more word for telling what happened.
\end{grammarlens}

\subsection*{Your turn}
\begin{itemize}
\raggedright
  \item \emph{Say it:} \emph{zǎonián}
  \item \emph{Say it:} \emph{zǎonián}, once more
  \item \emph{Say it:} \emph{tā zǎonián zài wàiguó gōngzuò guo} — In his early years he worked abroad.
  \item \emph{Recall:} say the Mandarin for a script for a language, then the Mandarin for in those years long ago, then say \emph{zǎonián} again
\end{itemize}

\begin{tcolorbox}[breakable,colback=teal!4,colframe=teal!35!black,title={Before you move on}]
What does \zh{早年} mean? (\textquotedblleft{}One's early years; long ago\textquotedblright{}.) Say it once more.
\end{tcolorbox}
\section[\texorpdfstring{\zh{晚年} (wǎnnián)}{晚年 (wǎnnián)}]{\zh{晚年} (wǎnnián) — one's later years; old age}

\label{lesson:ZH-C322-wannian}



\begin{quote}
Before the new one: say the Mandarin for in those years long ago, then the Mandarin for in one's early years.
\end{quote}

\subsection*{You'll want to know: \zh{晚年}}
\textbf{\zh{晚年}} — \emph{wǎnnián} — \textquotedblleft{}one's later years; old age\textquotedblright{}.

\begin{quote}
\textbf{\zh{他晚年的时候回了老家}} \emph{tā wǎnnián de shíhou huí le lǎojiā} — In his later years he went back to his home town.
\end{quote}

\begin{cousinweb}
\textbf{\zh{晚年}} (\emph{wǎnnián}) is built from \zh{晚} 'late' and \zh{年} 'a year', both inherited from Middle Chinese, \emph{mjonX} and \emph{nen} (Guangyun, in Baxter's transcription), which glosses the first 'dusk'.
\end{cousinweb}

\begin{grammarlens}[title={What you've built}]
One more word for telling what happened.
\end{grammarlens}

\subsection*{Your turn}
\begin{itemize}
\raggedright
  \item \emph{Say it:} \emph{wǎnnián}
  \item \emph{Say it:} \emph{wǎnnián}, once more
  \item \emph{Say it:} \emph{tā wǎnnián de shíhou huí le lǎojiā} — In his later years he went back to his home town.
  \item \emph{Recall:} say the Mandarin for in those years long ago, then the Mandarin for in one's early years, then say \emph{wǎnnián} again
\end{itemize}

\begin{tcolorbox}[breakable,colback=teal!4,colframe=teal!35!black,title={Before you move on}]
What does \zh{晚年} mean? (\textquotedblleft{}One's later years; old age\textquotedblright{}.) Say it once more.
\end{tcolorbox}
\section[\texorpdfstring{\zh{一生} (yìshēng)}{一生 (yìshēng)}]{\zh{一生} (yìshēng) — all one's life; a lifetime}

\label{lesson:ZH-C322-yisheng}



\begin{quote}
Before the new one: say the Mandarin for in one's early years, then the Mandarin for in one's later years.
\end{quote}

\subsection*{You'll want to know: \zh{一生}}
\textbf{\zh{一生}} — \emph{yìshēng} — \textquotedblleft{}all one's life; a lifetime\textquotedblright{}.

\begin{quote}
\textbf{\zh{他一生都在这个学校工作}} \emph{tā yìshēng dōu zài zhège xuéxiào gōngzuò} — He worked at this school all his life.
\end{quote}

A tone change: \zh{一} is said with a falling tone, yì, before the level tone of shēng.

\begin{cousinweb}
\textbf{\zh{一生}} (\emph{yìshēng}) is built from \zh{一} 'one' and \zh{生} 'to live', both inherited from Middle Chinese, \emph{'jit} and \emph{srjaeng} (Guangyun, in Baxter's transcription), which calls the first the start of the numbers and glosses the second 'to grow'.
\end{cousinweb}

\begin{grammarlens}[title={What you've built}]
One more word for telling what happened.
\end{grammarlens}

\subsection*{Your turn}
\begin{itemize}
\raggedright
  \item \emph{Say it:} \emph{yìshēng}
  \item \emph{Say it:} \emph{yìshēng}, once more
  \item \emph{Say it:} \emph{tā yìshēng dōu zài zhège xuéxiào gōngzuò} — He worked at this school all his life.
  \item \emph{Recall:} say the Mandarin for in one's early years, then the Mandarin for in one's later years, then say \emph{yìshēng} again
\end{itemize}

\begin{tcolorbox}[breakable,colback=teal!4,colframe=teal!35!black,title={Before you move on}]
What does \zh{一生} mean? (\textquotedblleft{}All one's life; a lifetime\textquotedblright{}.) Say it once more.
\end{tcolorbox}
\section[huáiniàn]{huáiniàn — to cherish the memory of; to look back fondly on}

\label{lesson:ZH-C322-huainian}



\begin{quote}
Before the new one: say the Mandarin for in one's later years, then the Mandarin for all one's life.
\end{quote}

\subsection*{You'll want to know: huáiniàn}
\emph{huáiniàn} — \textquotedblleft{}to cherish the memory of; to look back fondly on\textquotedblright{}.

\begin{quote}
\emph{wǒ hěn huáiniàn xiǎoshíhou de shēnghuó} — I really miss my childhood days.
\end{quote}

Its characters come later; for now it is a word for the ear.

\begin{cousinweb}
\emph{huáiniàn} is built from \emph{huái} and \emph{niàn}, both inherited from Middle Chinese, \emph{hweaj} and \emph{nemH} (Guangyun, in Baxter's transcription), which glosses the first 'to hold to the breast' and 'to think of', and the second 'to think of'.
\end{cousinweb}

\begin{grammarlens}[title={What you've built}]
One more word for telling what happened.
\end{grammarlens}

\subsection*{Your turn}
\begin{itemize}
\raggedright
  \item \emph{Say it:} \emph{huáiniàn}
  \item \emph{Say it:} \emph{huáiniàn}, once more
  \item \emph{Say it:} \emph{wǒ hěn huáiniàn xiǎoshíhou de shēnghuó} — I really miss my childhood days.
  \item \emph{Recall:} say the Mandarin for in one's later years, then the Mandarin for all one's life, then say \emph{huáiniàn} again
\end{itemize}

\begin{tcolorbox}[breakable,colback=teal!4,colframe=teal!35!black,title={Before you move on}]
What does huáiniàn mean? (\textquotedblleft{}To cherish the memory of; to look back fondly on\textquotedblright{}.) Say it once more.
\end{tcolorbox}
